\section{Discussion}
The results identify where inexpensive extraction is most immediately useful. It can generate a large, source-linked candidate collection; support targeted review by topic, horizon, and scope; and enable explicit sensitivity analysis. Those capabilities are substantive even when extraction is imperfect and predictive superiority is unproven. The key restriction is that the candidate corpus should not be described as a uniformly validated set of historical forecasts.

Cost and validity are separate dimensions of that judgment. The LOC experiment shows that low projected API expenditure need not imply zero extraction capability. It also shows a real error tradeoff: the economical self-verifying configuration has more unmatched claims and fewer reference matches than the higher-recall configuration. The ProQuest human review then establishes that inclusion and field agreement require further attention in the actual economic domain. Because the experiments differ in source, task, and reference construction, we cannot combine their cost and quality numbers into a single cost-per-valid-claim estimate.

The central measurement result is stronger than a generic warning about noisy labels. On a fixed historical cohort, observed accuracy, baseline accuracy, and their difference respond differently to the scoring rule. Switching the employment benchmark changes the sign of the estimated relative advantage without dramatically changing the newspaper hit rate. This illustrates why a paper about forecast accuracy must state both the proposition being judged and the comparison forecast. A favorable raw percentage is insufficient evidence of skill.

Several limitations remain consequential. Retrieval coverage and missing processing denominators restrict representativeness. The positive-candidate audit leaves recall unknown. Two human reviews are now available, but inter-reviewer agreement and adjudication remain unmeasured; partial machine-label exposure and coarse horizon categories also limit conclusions about annotation reliability. Revised outcome data and modern pretrained extraction limit claims of historical real-time feasibility. Block intervals have few independent units and omit extraction and model-selection uncertainty. The separate modern extension adds evidence about transfer but includes prior-exposure uncertainty, uneven calibration, and a pronounced employment failure (Appendix~\ref{app:modern}).

The next study should address these limits directly. The existing two-reviewer candidate audit should be extended to a shared sample of complete documents, including documents without forecasts, before comparing inexpensive and expensive extractors. That design would measure recall, inclusion precision, field agreement, and adjudication uncertainty on common support. A call-level usage ledger and review-time record would permit full cost accounting. For economic validation, source-supported targets and endpoints should be fixed before inspecting realized correctness, followed by matched-baseline estimates for the audited population. Any broader error adjustment would require a documented sampling design and assumptions about how reviewed and unreviewed records relate.

Prospective evaluation should freeze preprocessing, target definitions, model candidates, and forecast publication times before outcomes mature. Such an exercise would test performance under a genuinely new sequence of economic events. The present results do not justify a causal effect of political polarization, the supremacy of recurrent networks, or a successful 2026 recession forecast. They justify a research program in which inexpensive extraction expands the observable historical record while source validation and explicit economic definitions constrain the conclusions drawn from it.

